\documentclass[11pt,a4paper]{article}

\usepackage[margin=1in]{geometry}
\usepackage[T1]{fontenc}
\usepackage[utf8]{inputenc}
\usepackage{lmodern}
\usepackage{setspace}
\usepackage{hyperref}
\usepackage{longtable}
\usepackage{array}
\usepackage{enumitem}
\usepackage{titlesec}
\usepackage{xcolor}

\hypersetup{
  colorlinks=true,
  linkcolor=blue,
  urlcolor=blue,
  pdftitle={LOCK-IN / Arsii Project Technical Documentation},
  pdfauthor={Project Team}
}

\setstretch{1.08}
\setlist[itemize]{leftmargin=1.8em}
\setlist[enumerate]{leftmargin=2.0em}
\titleformat{\section}{\large\bfseries}{\thesection}{0.7em}{}
\titleformat{\subsection}{\normalsize\bfseries}{\thesubsection}{0.6em}{}
\titleformat{\subsubsection}{\normalsize\itshape}{\thesubsubsection}{0.5em}{}

\title{\textbf{LOCK-IN (Arsii)\\Detailed Technical Project Documentation}}
\author{Flutter + Firebase + AI Tutoring Platform}
\date{April 19, 2026}

\begin{document}
\maketitle
\tableofcontents
\newpage

\section{Executive Summary}
LOCK-IN (package name: \texttt{arsii}) is a Flutter web-first educational platform focused on Informatics learning. It combines:
\begin{itemize}
  \item Firebase Authentication for account lifecycle,
  \item Cloud Firestore for profile, lesson, attempt, and analytics persistence,
  \item AI-powered workflows (OpenRouter API) for personalized onboarding, adaptive exercises, and tutor interactions,
  \item A dark-theme, modern UI with responsive desktop/mobile layouts.
\end{itemize}

The core value proposition is adaptive learning: each student receives custom recommendations and support based on diagnostic assessment outcomes, weak-topic signals, and ongoing exercise performance.

\section{Product Scope and Learning Model}
\subsection{Learning Domains}
The project is designed around five Informatics domains represented in onboarding/mock subject definitions and lesson topic IDs:
\begin{enumerate}
  \item Programming Basics
  \item Algorithms
  \item Data Structures
  \item Databases
  \item Web Basics
\end{enumerate}

\subsection{Pedagogical Strategy}
The app follows a staged flow:
\begin{enumerate}
  \item Authentication and account entry.
  \item Subject enrollment by the student.
  \item AI-generated diagnostic evaluation quiz.
  \item Profile initialization in Firestore (level + per-subject scores + onboarding completion).
  \item Continuous adaptive delivery of lessons/exercises + progress analytics + tutor assistance.
\end{enumerate}

This creates a loop where data from student interactions continuously refines next recommendations.

\section{Technology Stack}
\subsection{Runtime and Language}
\begin{itemize}
  \item Flutter SDK (Dart $\geq$ 3.7.0 per \texttt{pubspec.yaml})
  \item Material 3 UI system
\end{itemize}

\subsection{Primary Dependencies}
\begin{itemize}
  \item \texttt{firebase\_core}: Firebase initialization
  \item \texttt{firebase\_auth}: user authentication, verification, password reset
  \item \texttt{cloud\_firestore}: persistent app data and analytics storage
  \item \texttt{http}: OpenRouter HTTP calls
  \item \texttt{flutter\_dotenv}: secret loading from \texttt{assets/env/app.env}
  \item \texttt{google\_fonts}: typographic system (Inter)
\end{itemize}

\subsection{Build and Platform}
\begin{itemize}
  \item Configured Firebase options currently expose a \textbf{web} target in \texttt{firebase\_options.dart}.
  \item \texttt{web/index.html} defines a branded loading screen and metadata.
  \item \texttt{firebase.json} contains FlutterFire project configuration metadata.
\end{itemize}

\section{Repository and Code Organization}
\subsection{Top-Level Layout}
\begin{itemize}
  \item \texttt{lib/main.dart}: app bootstrap + auth gate routing.
  \item \texttt{lib/core}: shared config, services, theming, mock data.
  \item \texttt{lib/features}: domain modules (auth, onboarding, dashboard, lessons, exercises, progress, profile, welcome).
  \item \texttt{lib/shared/widgets}: reusable UI components and visual primitives.
  \item \texttt{assets/images}, \texttt{assets/env}: branding and environment files.
\end{itemize}

\subsection{Architectural Style}
The app follows a pragmatic feature-first layered architecture:
\begin{itemize}
  \item \textbf{UI Layer}: Screens and widgets by feature.
  \item \textbf{Service Layer}: Firebase and AI orchestration in \texttt{core/services} and feature services.
  \item \textbf{Model Layer}: DTO/domain objects under feature model folders.
  \item \textbf{Infrastructure Layer}: Firebase init, dotenv config, and network integrations.
\end{itemize}

There is no heavy state management framework; state is mostly managed with StatefulWidgets, local state variables, FutureBuilder/StreamBuilder, and service calls.

\section{Application Bootstrap and Global Flow}
\subsection{Startup Sequence (\texttt{main.dart})}
At launch:
\begin{enumerate}
  \item Flutter bindings are initialized.
  \item Dotenv attempts to load \texttt{assets/env/app.env}; failure is tolerated.
  \item Firebase initializes using \texttt{DefaultFirebaseOptions.currentPlatform}.
  \item \texttt{ArsiiApp} starts with \texttt{AppTheme.darkTheme}.
\end{enumerate}

\subsection{Auth Gate Routing Logic}
\texttt{AuthGate} listens to Firebase auth changes and routes users:
\begin{itemize}
  \item No authenticated user: \texttt{WelcomeScreen}
  \item Authenticated but onboarding not complete: \texttt{EnrollmentScreen}
  \item Authenticated and onboarding complete: \texttt{DashboardScreen}
\end{itemize}

Onboarding completion is verified through \texttt{StudentService().isOnboardingComplete()} with timeout fallback.

\section{Authentication and Account Services}
\subsection{Auth Service Responsibilities}
\texttt{core/services/auth\_service.dart} centralizes FirebaseAuth interactions:
\begin{itemize}
  \item \texttt{authStateChanges}
  \item \texttt{currentUser}
  \item \texttt{signUp(email, password)}
  \item \texttt{signIn(email, password)}
  \item \texttt{resetPassword(email)}
  \item \texttt{sendEmailVerification()}
  \item \texttt{reloadCurrentUser()}
  \item \texttt{signOut()}
\end{itemize}

\subsection{UX Integration}
Authentication actions are surfaced in:
\begin{itemize}
  \item \texttt{WelcomeScreen}: login and forgot-password handling
  \item \texttt{signup\_screen.dart}: account creation flow
  \item \texttt{DashboardScreen}: verification email send/refresh dialogs
  \item \texttt{ProfileScreen}: sign-out capability
\end{itemize}

Error messaging maps Firebase auth codes to user-friendly text in dashboard/account-related flows.

\section{Onboarding Pipeline}
\subsection{Step 1: Enrollment}
\texttt{EnrollmentScreen} lets students choose desired subjects from data modeled in \texttt{MockData.subjects}. The UI supports desktop/mobile layouts and captures selected subject IDs.

\subsection{Step 2: AI Evaluation}
\texttt{EvaluationScreen} orchestrates:
\begin{enumerate}
  \item quiz generation through \texttt{AiService.generateQuiz(subjectNames)},
  \item per-question answer tracking and instant feedback animation,
  \item final evaluation via \texttt{AiService.evaluateResults(...)},
  \item persistence of onboarding results using \texttt{StudentService.saveOnboardingResults(...)}.
\end{enumerate}

\subsection{Evaluation Outcome}
The resulting profile baseline includes:
\begin{itemize}
  \item inferred level (Beginner/Intermediate/Advanced),
  \item subject performance map,
  \item onboarding completion flag.
\end{itemize}

Navigation then resets stack and transitions to dashboard.

\section{Core Domain Model and Profile Schema}
\subsection{Student Profile Model}
\texttt{StudentProfile} encapsulates user-learning state, including:
\begin{itemize}
  \item identity fields (uid, name, email),
  \item level,
  \item selected subjects and progress summaries,
  \item weak topics,
  \item recommendations,
  \item current lesson context (subject/title/description/progress),
  \item streak and activity metadata,
  \item \texttt{onboardingComplete}.
\end{itemize}

Additional nested models include \texttt{StudentSubject}, \texttt{ActivityItem}, \texttt{WeakTopic}, and \texttt{Recommendation}. Objects support map serialization for Firestore writes.

\subsection{Student Service Functional Scope}
\texttt{StudentService} is the main abstraction over student document operations:
\begin{itemize}
  \item profile retrieval,
  \item onboarding completion checks,
  \item onboarding result persistence,
  \item progress and subject updates,
  \item current lesson updates.
\end{itemize}

This service is consumed by dashboard, lessons, exercises, tutor, and progress components.

\section{AI Integration Design}
\subsection{Configuration}
\texttt{ApiConfig} reads \texttt{OPENROUTER\_API\_KEY} from dotenv and sets endpoint/model constants:
\begin{itemize}
  \item URL: \texttt{https://openrouter.ai/api/v1/chat/completions}
  \item model: \texttt{nvidia/nemotron-3-super-120b-a12b:free}
\end{itemize}

\subsection{AiService Responsibilities}
\texttt{AiService} supports multiple adaptive functions:
\begin{itemize}
  \item quiz generation for onboarding,
  \item result evaluation and level assignment,
  \item lesson simplification text generation,
  \item additional example generation,
  \item adaptive exercise generation,
  \item short-answer evaluation,
  \item hint generation,
  \item tutor-reply generation.
\end{itemize}

\subsection{Reliability Behaviors}
AI HTTP layer includes:
\begin{itemize}
  \item timeout handling,
  \item retry loops,
  \item HTTP 429 backoff logic using \texttt{retry-after} when available,
  \item null/fallback behavior when key is unavailable or request fails.
\end{itemize}

This graceful-degradation pattern allows the app to remain usable even when AI calls fail.

\section{Lessons Personalization Engine}
\subsection{Service Source Strategy}
\texttt{FirestoreLessonsService.getPersonalizedLessons()} composes lesson views from:
\begin{enumerate}
  \item student profile selected topic IDs,
  \item Firestore topic lesson collections (ordered by \texttt{order}),
  \item fallback seed data from \texttt{lesson\_seed\_data.dart},
  \item per-user lesson progress records.
\end{enumerate}

\subsection{Personalized Buckets}
It generates three principal sets:
\begin{itemize}
  \item \textbf{Recommended}: score-ranked for level fit and weak-skill relevance.
  \item \textbf{Continue Learning}: in-progress lessons sorted by recent open time.
  \item \textbf{Review Weak Areas}: unresolved lessons linked to weak topics.
\end{itemize}

\subsection{Lesson Open Tracking}
\texttt{markLessonOpened(...)} updates:
\begin{itemize}
  \item student \texttt{lesson\_progress} document,
  \item current lesson context on student profile via \texttt{StudentService.updateCurrentLesson(...)}.
\end{itemize}

\section{Exercises and Assessment Engine}
\subsection{Exercise Acquisition}
\texttt{ExerciseService.loadExercisesForLesson(...)} follows this chain:
\begin{enumerate}
  \item load Firestore exercises,
  \item if insufficient, generate AI adaptive exercises using lesson context, weak skills, and mistakes,
  \item if generation fails/empty, return fallback template exercises.
\end{enumerate}

\subsection{Answer Evaluation}
\begin{itemize}
  \item Multiple choice: deterministic comparison to expected answer.
  \item Short answer: AI-based semantic evaluation returning correctness, score, explanation, and personalized feedback.
\end{itemize}

\subsection{Persistence}
Each exercise interaction can store:
\begin{itemize}
  \item individual attempt in \texttt{exercise\_attempts},
  \item session metadata in \texttt{exercise\_sessions},
  \item final session result with weak-area extraction,
  \item downstream updates to subject mastery and lesson progress status.
\end{itemize}

\section{Progress Analytics Pipeline}
\subsection{Computation Source}
\texttt{FirestoreProgressService.getSummary()} aggregates:
\begin{itemize}
  \item student profile baseline,
  \item lesson progress collection,
  \item latest exercise attempts (bounded query).
\end{itemize}

\subsection{Derived Metrics}
Computed metrics include:
\begin{itemize}
  \item lesson completion ratio,
  \item exercises attempted,
  \item accuracy rate,
  \item mastery by skill,
  \item weak areas by mistake concentration,
  \item trend string comparing recent vs older attempt windows,
  \item recommendation text for next review action.
\end{itemize}

\subsection{Output and Storage}
Summary is returned to UI and merged into \texttt{students/{uid}/progress\_summary/main}, supporting quick retrieval and historical updates.

\section{AI Tutor Conversation Subsystem}
\subsection{Message Flow}
\texttt{AiTutorService.askTutor(...)}:
\begin{enumerate}
  \item loads profile context (level, weak topics, progress),
  \item loads recent incorrect exercise attempts,
  \item invokes AI tutor reply generation,
  \item writes user and tutor messages into \texttt{tutor\_messages} subcollection in one transaction,
  \item returns tutor message DTO for rendering.
\end{enumerate}

\subsection{Safety and Query Design}
To reduce Firestore index complexity, some filtering is performed client-side after ordered queries (notably lesson/message filtering and incorrect-attempt extraction patterns).

\section{Feature-by-Feature UI Documentation}
\subsection{Welcome and Auth Entry}
\begin{itemize}
  \item Animated split-screen desktop and stacked mobile layout.
  \item Branded top navigation with informational dialogs (features, curriculum, about).
  \item Login form with email/password validation and mapped error states.
\end{itemize}

\subsection{Onboarding}
\begin{itemize}
  \item Enrollment step with selectable cards.
  \item AI evaluation step with progress bar, per-question interaction, and result summary.
\end{itemize}

\subsection{Dashboard}
\begin{itemize}
  \item Serves as post-onboarding hub and tabbed shell (home, lessons, progress, exercises, settings).
  \item Integrates profile loading, email verification actions, and navigation into sub-features.
\end{itemize}

\subsection{Lessons}
\begin{itemize}
  \item Categorized personalized lists.
  \item Adaptive cards showing reasons and badges.
  \item Lesson details page with content, concepts, examples, and AI action buttons.
  \item Embedded/side AI tutor panel depending on viewport width.
\end{itemize}

\subsection{Exercises}
\begin{itemize}
  \item Session-based question workflow per lesson.
  \item Multi-type answer input support.
  \item Real-time feedback, hints, and final summary.
\end{itemize}

\subsection{Progress}
\begin{itemize}
  \item Hero analytics card and KPI cards.
  \item Insight and mastery breakdown sections.
  \item Empty/error states and refresh behaviors.
\end{itemize}

\subsection{Profile}
\begin{itemize}
  \item Identity card with account metadata.
  \item Sign-out flow returning to welcome/login context.
\end{itemize}

\section{Theme, Design Tokens, and Shared Components}
\subsection{Color System}
\texttt{AppColors} defines a dark, neon-accent palette:
\begin{itemize}
  \item background surfaces (dark/card/elevated/subtle),
  \item primary cyan spectrum,
  \item secondary violet and tertiary pink accents,
  \item semantic status colors,
  \item reusable gradients and glow helpers.
\end{itemize}

\subsection{Typography and Material Theme}
\texttt{AppTheme.darkTheme} applies Material 3 + Inter text theme with custom hierarchy, button styles, card borders, and dark ColorScheme integration.

\subsection{Shared Widget Library}
Reusable widgets in \texttt{lib/shared/widgets} provide consistent visuals and interactions, including:
\begin{itemize}
  \item abstract animated background,
  \item glow/styled cards,
  \item gradient button,
  \item styled text field,
  \item dashboard cards and hero components,
  \item onboarding option and assessment card components.
\end{itemize}

\section{Data and Firestore Collection Model}
\subsection{Main Collection Topology}
A practical inferred schema from service code:
\begin{itemize}
  \item \texttt{students/{uid}}
  \begin{itemize}
    \item profile fields (level, subjects, weakTopics, recommendations, current lesson context, streak, onboardingComplete, ...)
    \item \texttt{lesson\_progress/{lessonId}}
    \item \texttt{exercise\_attempts/{attemptId}}
    \item \texttt{exercise\_sessions/{sessionId}}
    \item \texttt{progress\_summary/main}
    \item \texttt{tutor\_messages/{messageId}}
  \end{itemize}
  \item \texttt{topics/{topicId}/lessons/{lessonId}}
\end{itemize}

\subsection{Normalization Strategy}
The app mixes normalized and denormalized structures:
\begin{itemize}
  \item normalized lesson repository by topic,
  \item denormalized user analytics snapshots (progress summary),
  \item append-only attempts/messages for history and trend analysis.
\end{itemize}

\section{Environment and Secrets Handling}
\subsection{Dotenv Use}
AI key retrieval uses \texttt{flutter\_dotenv} from \texttt{assets/env/app.env}. If missing/invalid, AI functions gracefully degrade.

\subsection{Security Consideration}
Because the current app is web-oriented and calls external AI APIs, production hardening should eventually move sensitive key usage behind a secure backend proxy to avoid exposing provider secrets in client-delivered contexts.

\section{Error Handling and Resilience Patterns}
Project-wide patterns include:
\begin{itemize}
  \item guarding operations with mounted checks in async UI code,
  \item explicit loading/error/empty states per screen,
  \item timeouts on network-critical onboarding operations,
  \item tolerant fallbacks for AI and content sourcing,
  \item retry actions in UI for recoverable failures.
\end{itemize}

\section{Responsive and UX Considerations}
\begin{itemize}
  \item Repeated desktop breakpoints around wider layouts ($\sim$960 px and above).
  \item Side-panel tutor for very wide screens and floating compact panel on smaller layouts.
  \item Visual hierarchy supports high contrast in dark mode with accent-guided focus.
\end{itemize}

\section{Current Operational Constraints and Risks}
\subsection{Platform Scope}
\texttt{firebase\_options.dart} currently throws unsupported error for non-web platforms. If mobile/desktop deployment is required, additional platform configs must be generated and integrated.

\subsection{Client-Side AI Calls}
Direct API calls from client-side runtime can introduce key-management and abuse risks. Recommended future direction: server-side token mediation or Cloud Functions gateway.

\subsection{Firestore Index and Query Strategy}
Some client-side filtering choices reduce immediate index requirements but may increase over-fetching. This trade-off is acceptable for early-stage scale but should be revisited with production load.

\section{Build, Run, and Maintenance Notes}
\subsection{Typical Local Setup}
\begin{enumerate}
  \item Ensure Flutter and Dart SDK compatibility with \texttt{pubspec.yaml}.
  \item Provide \texttt{assets/env/app.env} with \texttt{OPENROUTER\_API\_KEY}.
  \item Configure Firebase project credentials (already scaffolded for web).
  \item Run dependency install and launch Flutter web target.
\end{enumerate}

\subsection{Testing Status}
The repository includes Flutter test dependency scaffolding, but this documentation reflects architectural review of implementation code and does not claim full automated test coverage of all flows.

\section{Suggested Future Enhancements}
\begin{enumerate}
  \item Introduce structured state management (e.g., Riverpod/Bloc) for larger scalability.
  \item Add repository interfaces and test doubles for easier unit/integration testing.
  \item Move AI provider calls to backend services for secure key handling.
  \item Add analytics event instrumentation for pedagogical outcome tracking.
  \item Expand offline/cache strategies for lesson/exercise continuity.
  \item Strengthen CI with lint, static analysis, widget tests, and end-to-end smoke tests.
\end{enumerate}

\section{Conclusion}
LOCK-IN (Arsii) is a substantial adaptive-learning Flutter application with clear feature modularization, rich UI implementation, and meaningful AI + Firebase integration. The current architecture is practical for iterative development and already contains key foundations of an intelligent tutoring platform: personalized diagnostics, adaptive content generation, learner-progress analytics, and context-aware tutor support.

With security hardening, expanded test coverage, and backend mediation for AI traffic, the project is well-positioned for production-level maturation.

\end{document}
